% Per-game replacements of the fixed texts of lang-en.tex (\renewcommand).
\renewcommand{\tCoverLine}{The MSX turbo R game on the V9990}
\renewcommand{\tMemoryText}{\textbf{\ramNeeded{} of RAM}; every MSX turbo R has it.}
\renewcommand{\tVideoText}{A \textbf{V9990 cartridge} (Yamaha V9990 video). It is mandatory: the whole picture is drawn by the V9990. Apart from the messages at power-on, the computer's own screen is not used by the game.}
\renewcommand{\tSoundText}{Connect your speakers to the \textbf{MoonSound or Wozblaster} to hear music and effects together. The cartridge mixes FM music with PCM gunfire, explosions and other arcade sounds. With an OPL3-only cartridge the MSX audio output carries the PSG effects and the cartridge supplies the music. Without a sound cartridge, use the MSX output for effects.}
\renewcommand{\tSetGame}{Insert the \textbf{\gameTitle{} cartridge} in another slot: a primary slot or a subslot of a slot expander. The order does not matter.}
\renewcommand{\tSetOn}{Switch the MSX on. For a moment the MSX's own screen shows \texttt{V9990 found !}; then the picture appears on the V9990's monitor. If it says \texttt{THIS GAME RUNS ONLY ON MSX TURBO R} the computer is not a turbo R, and if it says \texttt{V9990 NOT found!!!} the V9990 is missing or not seated; in both cases the game stops there.}
\renewcommand{\tControlsText}{Keyboard (cursor keys), or a joystick with two buttons in port 1.}
\renewcommand{\tPowerOnCaption}{Left: the cover. Right: the title screen with the best seven scores.}
\renewcommand{\tCreditsTitle}{10. Credits and rights}
\newcommand{\tCreditsExtra}{\item The cover picture is the artwork of the original game's home versions. The history in section~5 is a summary of public sources (Wikipedia and others).}
